\documentclass{ximera}
%geschreven 13/11/26
\title{Behoud van energie}
\author{Daan Lipkens}

\begin{document}
\begin{abstract}
    De wet van behoud van energie, de arbeid-energiestelling en het verband tussen arbeid en potentiële energie.
\end{abstract}
\maketitle

\section{De wet van behoud van energie}\label{sec:behoud_energie}
Energie kan niet ontstaan of verdwijnen: ze wordt enkel omgezet van de ene vorm in de andere. Laat je een steen vallen, dan daalt zijn potentiële energie en stijgt zijn kinetische energie (figuur~\ref{fig:vallende_steen}).

\begin{figure}[h]
    \captionsetup{width=0.85\textwidth}
    \begin{center}
        \begin{tikzpicture}
            % grond
            \draw[black!60, line width=0.8pt] (-1.5,0) -- (6.5,0);
            \foreach \x in {-1.4,-1.1,...,6.4} \draw[black!40] (\x,0) -- (\x-0.2,-0.2);
            \node[black!60, below] at (-0.9,-0.2) {RP};
            % hoogte-as
            \draw[-{Latex[length=2mm]}, draw=black!50, text=black!50] (-0.9,0) -- (-0.9,4.4) node[above] {$h$};
            % drie standen
            \foreach \y/\lab in {4/1, 2/2, 0.25/3} {
                \fill[black!70] (0,\y) circle (0.22);
                \node[right] at (1.1,\y) {\lab};
            }
            \draw[snelheid] (0.5,2.25) -- (0.5,1.75) node[midway, right] {$\vec{\mathbf{v}}$};
            \draw[snelheid] (0.5,0.7) -- (0.5,0.0) node[midway, right] {$\vec{\mathbf{v}}$};
            % energiebalkjes
            \foreach \y/\p/\k in {4/1.6/0, 2/0.8/0.8, 0.25/0/1.6} {
                \draw[black!50] (2,\y-0.2) rectangle (3.6,\y+0.2);
                \fill[blue!50] (2,\y-0.2) rectangle (2+\p,\y+0.2);
                \fill[red!50] (2+\p,\y-0.2) rectangle (3.6,\y+0.2);
            }
            \fill[blue!50] (4.2,3.2) rectangle (4.5,3.45); \node[right] at (4.5,3.32) {$E_{\text{pot}}$};
            \fill[red!50] (4.2,2.7) rectangle (4.5,2.95); \node[right] at (4.5,2.82) {$E_{\text{kin}}$};
        \end{tikzpicture}
    \end{center}
    \caption{Een steen valt vanuit rust (stand 1) tot vlak boven de grond (stand 3). Zonder luchtwrijving blijft de som van potentiële en kinetische energie gelijk: de balkjes hebben telkens dezelfde lengte.}
    \label{fig:vallende_steen}
\end{figure}

\begin{proposition}[title={Wet van behoud van mechanische energie}]\nl
    Werken er geen wrijvingskrachten, dan blijft de som van de potentiële en de kinetische energie constant. Voor twee standen $1$ en $2$ geldt:
    \begin{equation}\label{eq:behoud_energie}
        E_{\text{pot}}(1) + E_{\text{kin}}(1) = E_{\text{pot}}(2) + E_{\text{kin}}(2)
    \end{equation}
\end{proposition}

\begin{remark}[title={Met wrijving}]\nl
    Werkt er wel wrijving, dan wordt een deel van de mechanische energie omgezet in warmte. De totale energie blijft behouden, maar de som $E_{\text{pot}} + E_{\text{kin}}$ daalt. Zo botsen de vrije elektronen in een geleider tegen de roosterionen en wordt hun kinetische energie omgezet in warmte.
\end{remark}

\begin{example}[foldable=true,title={Voorbeeld: een vallende steen}]\nl
    Een steen wordt vanuit rust losgelaten op $5{,}0 \, \mathrm{m}$ boven de grond. Bereken de snelheid waarmee hij de grond raakt. Verwaarloos de luchtwrijving.

    \begin{solution}\nl
        \underline{Gegevens:}
        \begin{align*}
            & h_{1} = 5{,}0 \, \mathrm{m} & & v_{1} = 0 \, \mathrm{m/s} & & h_{2} = 0 \, \mathrm{m} \text{ (grond = RP)}
        \end{align*}

        \underline{Gevraagd:} $v_{2}$

        \underline{Oplossing:}\\
        Er werkt geen wrijving, we gebruiken dus de wet van behoud van energie (vergelijking~\ref{eq:behoud_energie}). In stand 1 is $E_{\text{kin}} = 0 \, \mathrm{J}$, in stand 2 is $E_{\text{pot}} = 0 \, \mathrm{J}$.
        \begin{align*}
            m \cdot g \cdot h_{1} &= \frac{1}{2} \cdot m \cdot v_{2}^{2} && \text{(behoud van energie, $m$ valt weg)}\\
            v_{2} &= \sqrt{2 \cdot g \cdot h_{1}} && \text{(omvormen naar $v_{2}$)}\\
            &= \sqrt{2 \cdot 9{,}81 \, \mathrm{m/s^2} \cdot 5{,}0 \, \mathrm{m}} && \text{(gegevens invullen)}\\
            &= 9{,}90 \, \mathrm{m/s} \simeq \boxed{9{,}9 \, \mathrm{m/s}} && \text{(uitrekenen)}
        \end{align*}
    \end{solution}
\end{example}

\section{Arbeid en energie}\label{sec:arbeid_energie}
Arbeid is energie die door een kracht wordt overgedragen. Verricht een kracht positieve arbeid op een voorwerp, dan krijgt het voorwerp energie. Bij negatieve arbeid verliest het voorwerp energie.

\begin{proposition}[title={Arbeid-energiestelling}]\nl
    De totale arbeid die de krachten op een voorwerp verrichten, is gelijk aan de verandering van zijn kinetische energie:
    \begin{equation}\label{eq:arbeid_energiestelling}
        W = \Delta E_{\text{kin}} = E_{\text{kin}}(2) - E_{\text{kin}}(1)
    \end{equation}
\end{proposition}

\begin{proposition}[title={Arbeid en potentiële energie}]\nl
    Verricht de zwaartekracht (of de elektrische kracht) een arbeid $W$, dan verandert de potentiële energie met de tegengestelde waarde:
    \begin{equation}\label{eq:arbeid_epot}
        \Delta E_{\text{pot}} = -W
    \end{equation}
    Valt een steen, dan verricht de zwaartekracht positieve arbeid en daalt de potentiële energie.
\end{proposition}

\begin{example}[foldable=true,title={Voorbeeld: een kar versnellen}]\nl
    Een kar van $2{,}0 \, \mathrm{kg}$ staat stil op een horizontale baan zonder wrijving. Je duwt de kar $3{,}0 \, \mathrm{m}$ vooruit met een constante horizontale kracht van $12 \, \mathrm{N}$. Bereken de eindsnelheid van de kar.

    \begin{solution}\nl
        \underline{Gegevens:}
        \begin{align*}
            & m = 2{,}0 \, \mathrm{kg} & & F = 12 \, \mathrm{N} & & \Delta x = 3{,}0 \, \mathrm{m} & & v_{1} = 0 \, \mathrm{m/s}
        \end{align*}

        \underline{Gevraagd:} $v_{2}$

        \underline{Oplossing:}\\
        Eerst berekenen we de arbeid van de duwkracht. Kracht en verplaatsing hebben dezelfde zin ($\alpha = 0^\circ$), zie vergelijking~\verwijs[Arbeid]{eq:arbeid}{appendix_energie_arbeid}.
        \begin{align*}
            W &= F \cdot \Delta x = 12 \, \mathrm{N} \cdot 3{,}0 \, \mathrm{m} = 36 \, \mathrm{J} && \text{(arbeid)}
        \end{align*}
        Die arbeid wordt volledig kinetische energie, want de kar start in rust (vergelijking~\ref{eq:arbeid_energiestelling}).
        \begin{align*}
            W &= \frac{1}{2} \cdot m \cdot v_{2}^{2} - 0 && \text{(arbeid-energiestelling)}\\
            v_{2} &= \sqrt{\frac{2 \cdot W}{m}} = \sqrt{\frac{2 \cdot 36 \, \mathrm{J}}{2{,}0 \, \mathrm{kg}}} && \text{(omvormen en invullen)}\\
            &= \boxed{6{,}0 \, \mathrm{m/s}} && \text{(uitrekenen)}
        \end{align*}
    \end{solution}
\end{example}

\begin{question}
    Een bal rolt een helling af. Er is geen wrijving. Welke uitspraken zijn juist?

    \begin{selectAll}
        \choice[correct]{De potentiële energie van de bal daalt.}
        \choice[correct]{De zwaartekracht verricht positieve arbeid op de bal.}
        \choice{De som van potentiële en kinetische energie stijgt.}
        \choice[correct]{De kinetische energie stijgt evenveel als de potentiële energie daalt.}
    \end{selectAll}

    \begin{solution}\nl
        Zonder wrijving blijft $E_{\text{pot}} + E_{\text{kin}}$ constant. De zwaartekracht heeft een component in de zin van de verplaatsing, verricht dus positieve arbeid, en daardoor daalt $E_{\text{pot}}$ ($\Delta E_{\text{pot}} = -W$) en stijgt $E_{\text{kin}}$ met dezelfde hoeveelheid.
    \end{solution}
\end{question}

\end{document}
